\section{The deduction and the status of the argument}\label{sec:conclusion}
We first isolate the precise logical role of the computational assertion.

\begin{theorem}[Certificate-to-game implication]\label{thm:bridge}
If the finite exclusion and input checks in
Proposition~\ref{prop:finite} are correct, then there is no five-color
winning strategy on $C_5$.
\end{theorem}
\begin{proof}
Suppose a winning strategy exists. Proposition~\ref{prop:coverage}
places an equivalent normalized winning strategy in a parent count branch
with $17\leq n\leq25$. If the parent is a direct terminal branch,
Proposition~\ref{prop:input} extends this strategy to a satisfying
assignment of $F_a$, whereas Proposition~\ref{prop:finite} and
Lemma~\ref{lem:rup} make $F_a$ unsatisfiable.

If the parent is refined, Lemma~\ref{lem:patterns} first places an
equivalent strategy in one of its 51 subcases. This transformation
preserves the parent counts and its normalization. The same extension
argument gives a satisfying assignment of that $F_{a,t}$, contradicting
its checked refutation. These two cases exhaust the parent list.
\end{proof}

\begin{proof}[Argument for Theorem~\ref{thm:main}]
The original verification and the complete new replay supply the premise of
Theorem~\ref{thm:bridge}, and hence excludes five colors, subject to the
verification status stated in this manuscript. Proposition~\ref{prop:four}
provides a winning strategy for four colors. If a winning strategy existed
for more than five colors, Lemma~\ref{lem:palette} would produce one for
five colors. Therefore the value established by this computer-assisted
argument is exactly four.
\end{proof}

There is also an informative intermediate consequence of excluding the
weights below~25.
\begin{corollary}[The last remaining count]\label{cor:balance}
Assume that the exclusions at weights 17--24 are correct. Any five-color
winning strategy would then satisfy
\[
S(1)=S(2)=S(3)=S(4)=S(5)=25
\]
and would have no entry that guesses a neighbor's color. The verified
exclusion at weight 25, if correct, contradicts this necessary condition.
\end{corollary}
\begin{proof}
For any fixed $z$, the count-preserving normalization of
Lemma~\ref{lem:normalization}, followed by the same graph and color
reductions, would place a strategy with $S(z)\leq24$ in an excluded branch;
Proposition~\ref{prop:seventeen} rules out totals below~17. Thus every
$S(z)\geq25$. Lemma~\ref{lem:total} gives $\sum_zS(z)\leq125$ even before
normalization. Equality must hold in all five lower bounds, and every
entry must contribute to the sum. An entry guessing a neighbor's color
would contribute zero, so no such entry can occur. Weight 25 is then
excluded by the final family of certificates.
\end{proof}

This equality concerns the total for a color over all five players; it
does not say that each player guesses each color five times. It is a
necessary condition inside the proposed contradiction, not a construction
of a winning strategy. Moreover, none of the branch inputs used another
color's lower bound or the conclusion of another weight. The finite
exclusions therefore do not depend circularly on this balance statement.

\subsection*{Computer and AI assistance}
OpenAI Codex assisted with mathematical exploration, program development,
local computations, organization of the evidence, and preparation and
revision of this manuscript. The solver, input auditors, and proof
checkers have different roles, but were integrated within the same
AI-assisted workflow. The status distinction made in the introduction
therefore applies throughout: separate local programs are not external
reviewers, and neither the written argument nor the checking software has
been formally verified in a proof assistant.

\subsection*{Data availability and further work}
The organized research archive contains definitions, proof drafts,
programs, inputs, certificates, audit reports, and historical failures.
A private repository copy and versioned evidence attachments exist for
the earlier manuscript and research archive. The complete audit package
and the present revised manuscript have been organized locally; their
creation is not a claim of public availability or external review.
The appendix gives paths relative to the research root, so readers with
access can trace each assertion to its source and reproduce the checks.
Manuscript source alone is insufficient for a full replay: the large
certificate files must also be present.

A full local replay with newly implemented software has now been
completed. The next verification tasks are external scrutiny of the
combinatorial reductions and use of a separately reviewed or formally
verified certificate checker. A further mathematical goal is to replace the large finite
exclusion with a shorter argument explaining the incompatibility between
the three occurrence families. The local witness shows why the
240 two-occurrence colorings alone cannot supply that explanation.
